\documentclass[a4paper]{article}
\usepackage[pdftex]{graphicx}
\usepackage[utf8]{inputenc}
\usepackage{enumerate}
\usepackage{colortbl}
\usepackage{siunitx}
\sisetup{locale=DE}
\usepackage{href-ul}
\hypersetup{
	colorlinks=true,
	linkcolor=blue,
	urlcolor=blue}
\usepackage{icomma}
\usepackage{amssymb}
\usepackage{geometry}
\geometry{a4paper, top=15mm, left=15mm, right=15mm, bottom=15mm,
	headsep=10mm, footskip=12mm}

\begin{document}
	
	\begin{enumerate}[1.]
		
		{\bf \item Structure of Atoms}
		
		\begin{minipage}{0.65\textwidth}
			{\bf Rutherford bombarded a very thin gold foil with $\alpha$ particles} (helium nuclei, doubly positively charged) from a radioactive preparation and observed whether and at what angle they were deflected. He got the following surprising results:
		\end{minipage}
		\hfill
		\begin{minipage}{0.3\textwidth}
			\includegraphics[width=4 cm]{Rutherford129.pdf}
		\end{minipage}
		
		\begin{itemize}
			\item most $\alpha$ particles pass through the film unhindered.
			\item A few are deflected at sometimes large angles, very few are even backscattered.
		\end{itemize}
		
		\begin{enumerate}[a)]
			\item Describe the structure of an atom and explain why the majority of it
			$\alpha$ particles penetrate the film undeflected.
			\item Explain the surprising phenomenon that some $\alpha$ particles are deflected or even backscattered. Address the forces at work.
		\end{enumerate}
		
		{\bf \item The conventional notation for the nucleus of a barium atom is ${}^{138}_{~56}{\mathbf{Ba}}$}
		\begin{enumerate}[a)]
			\item Give all the information you can glean from this notation.
			\item Explain why the barium atom is electrically neutral despite the positively charged nucleus.
		\end{enumerate}
		
		{\bf \item half-life}
		
		\begin{minipage}{0.6\textwidth}
			At time $t=0$, a rock sample contains $\SI{1,000}{\gram}$ of the thallium isotope ${}^{206}_{~81}{\mathbf{Tl}}$ and no lead. ${}^{206}_{~81}{\mathbf{Tl}}$ converts into the stable lead after a $\beta$ decay. The half-life for this is 4.2 minutes.
			\begin{enumerate}[a)]
				\item Write the nuclear reaction equation for the $\beta$ decay of thallium-206.
				\item Fill in the table.
				\item Is there a point in time when everything is completely transformed? Reasons!
			\end{enumerate}
		\end{minipage}
		\hspace{0.5 cm}
		\begin{minipage}{0.4\textwidth}
			\renewcommand{\arraystretch}{2}
			\begin{tabular}{|>{\columncolor[gray]{.8}}p{1.8 cm}|p{1.5 cm}|p{1.5 cm}|}
				\hline
				\rowcolor[gray]{.8} Time in min & mass of thallium & mass of lead \\
				\hline
				0 & $\SI{1000}{\milli\gram}$ & $\SI{0}{\milli\gram}$\\
				\hline
				4.2 & &\\
				\hline
				8.4& &\\
				\hline
				12.6& &\\
				\hline
				42& &\\
				\hline
			\end{tabular}
		\end{minipage}
		
		{\bf \item Protons and neutrons consist of so-called quarks.}\\ Describe the exact structure of the nucleons. Also consider the electrical charge.
		
		{\bf \item Name three properties} of $\alpha$--, $\beta$--, and $\gamma$--radiation
		
		{\bf \item Radon ${}^{218}_{~86}\mathbf{Rn}$ was created} from an unknown radioactive isotope after an $\alpha$ decay.\\ Write up the nuclear reaction equation.
			
			{\bf \item Strontium ${}^{90}_{38}\mathbf{Sr}$ with a half-life of 29 years} is a $\beta$ emitter and one of the most dangerous secondary products of above-ground nuclear weapons tests. After the tests were almost completely stopped, an Sr--90-- activity of $\SI{2.5}{\kilo\becquerel}$ per square meter was measured in Neuherberg near Munich in 1966.
			
			\begin{enumerate}[a)]
				\item Write the nuclear reaction equation for the $\beta$ decay of strontium-90.
				\item Calculate the activity expected per square meter in 2020.
				\item Determine the year from which 75\% decayed.
			\end{enumerate}
			
			%Task 4
			{\bf \item In the naturally occurring thorium decay series} the radioactive changes
			${}^{232}_{~90}\mathbf{Th}$ into the stable lead isotope in several steps
			${}^{208}_{~82}\mathbf{Pb}$ around.\\
			Determine the number of each occurring by calculation
			$\alpha$-- and $\beta$--decays.
			
		\end{enumerate}
		
		\fbox{
			\begin{minipage}{0.5\textwidth}
				For the solution, please \href{https://www.okuyakl.de/phys/p9atoL412/le412.pdf}{click here} or scan the QR code.\\
				Additional worksheets are available at
				
				\href{http://www.okuyakl.com}{www.okuyakl.com}
			\end{minipage}
			\hfill
			\begin{minipage}{0.4\textwidth}
				\includegraphics[width=1.5 cm]{../../viecher/zwe03}
				\includegraphics[width=3 cm]{qre412}
				\includegraphics[width=2 cm]{../../viecher/afanticon1}
				
		\end{minipage}}
		
	\end{document}